% Authors: Carmen Beltran, Brendan Butler, Teddy Wachtler
%%%%%%%%%%%%%%%%%%%%%%%%%%%%%%%%%%%%%%%%%%%%%%%

\documentclass{article}
\usepackage[dvips]{graphicx}
\usepackage{a4wide}
\usepackage{amsmath}
\usepackage{euscript}
\usepackage{amssymb}
\usepackage{amsthm}
\usepackage{amsopn}

\theoremstyle{definition}
\newtheorem*{definition}{Definition}
\newtheorem{theorem}{Theorem}

\renewcommand{\qed}{\blacksquare}
\newcommand{\powset}{{\cal P}}
\newcommand{\img}{_{*}}
\newcommand{\pimg}{^{*}}

\newcommand{\vv}{\ensuremath{\vec{v}}}
\newcommand{\vu}{\ensuremath{\vec{u}}}
\newcommand{\vw}{\ensuremath{\vec{w}}}
\newcommand{\vx}{\ensuremath{\vec{x}}}
\newcommand{\vy}{\ensuremath{\vec{y}}}
\newcommand{\vb}{\ensuremath{\vec{b}}}
\newcommand{\vo}{\ensuremath{\vec{0}}}
\newcommand{\va}{\ensuremath{\vec{a}}}
\newcommand{\ve}{\ensuremath{\vec{e}}}

\newcommand{\R}{\mathbb{R}}
\newcommand{\Z}{\mathbb{Z}}
\newcommand{\C}{\mathbb{C}}
\newcommand{\N}{\mathbb{N}}
\newcommand{\Q}{\mathbb{Q}}

%%%%%%%%%%%%%%%%%%%%%%%%%%%%%%%%%%%%%%%%%%%%%%%

\begin{document}
	\begin{center}
		\Large{Math 251W: Foundations of Advanced Mathematics}
		\vspace{0.1cm}

		\normalsize Portfolio Assignment 5: \S 5.1 \& 5.3
		\vspace{0.1cm}

		\hfill Name: Carmen Beltran, Brendan Butler, Teddy Wachtler
		\vspace{0.1cm}

		\hrule
		\vspace{0.5cm}
	\end{center}

%%%%%%%%%%%%%%%%%%%%%%%%%%%%%%%%%%%%%%%%%%%%%%%

	\noindent Problem 5.1.8
	\vspace{0.2cm}

	Let $A$ be a set.
	Let $\{R_i\}_{i \in I}$ be a family of relations on $A$ indexed by the nonempty set $I$.

	\vspace{0.2cm}
	\begin{enumerate}
		\item[a] Prove each of the following propositions:

		\begin{enumerate}
			\item[i]
			If $R_i$ is reflexive for all $i \in I$ then
			$\bigcap_{ i\in I} R_i$ is reflexive

			\vspace{0.2cm}
			\fbox{Proof} ( )
			\vspace{0.2cm}

			$\qed$.

			\vspace{0.2cm}

			\item[ii]
			If $R_i$ is symmetric for all $i \in I$ then
			$\bigcap_{i \in I} R_i$ is symmetric

			\vspace{0.2cm}
			\fbox{Proof} ( )
			\vspace{0.2cm}

			$\qed$.

			\vspace{0.2cm}

			\item[iii]
			If $R_i$ is transitive for all $ i\in I$ then
			$\bigcap_{i \in I} R_i$ is transitive

			\vspace{0.2cm}
			\fbox{Proof} ( )
			\vspace{0.2cm}

			$\qed$.

			\vspace{0.2cm}

		\end{enumerate}

		\vspace{0.2cm}

		\item[b] Find a counterexample to the proposition: If $R_i$ is transitive for all $i \in I$ then $\bigcup_{i \in I} R_i$ is transitive.

		\vspace{0.2cm}
		\fbox{Counterexample} ( )
		\vspace{0.2cm}

		$\qed$.

		\vspace{0.5cm}

	\end{enumerate}

%%%%%%%%%%%%%%%%%%%%%%%%%%%%%%%%%%%%%%%%%%%%%%%

	\noindent Problem 5.1.11
	\vspace{0.2cm}

	\begin{enumerate}
		\item[3]
		\underline{Proposition:}
		If ${\cal R}$ is a transitive relation on the set $A$ and $x{\cal R}y$ then $[y] \subseteq [x]$

	\end{enumerate}

	\vspace{0.2cm}
	\fbox{Proof} ( )
	\vspace{0.2cm}

	$\qed$.

	\vspace{0.5cm}

%%%%%%%%%%%%%%%%%%%%%%%%%%%%%%%%%%%%%%%%%%%%%%%

	\noindent Problem 5.3.6
	\vspace{0.2cm}

	Let $A$ be a non-empty set, let $\sim$ be an equivalence relation on $A$, and let $A / \sim$ be the set of equivalence classes (i.e. $A / \sim \quad = \{[a]|a \in A\}$). Note it is common to use $\sim$ instead of $\mathcal{R}$ when the relation is an equivalence relation.

	\vspace{0.2cm}
	\underline{Proposition:} Let $x, y \in A$.

	\begin{enumerate}
		\item[a]
		If $x \sim y$, then $[x] = [y]$

		\vspace{0.2cm}
		\fbox{Proof} ( )
		\vspace{0.2cm}

		$\qed$.

		\vspace{0.2cm}

		\item[b]
		If $x \not \sim y$ then $[x] \cap [y] = \emptyset$

		\vspace{0.2cm}
		\fbox{Proof} ( )
		\vspace{0.2cm}

		$\qed$.

		\vspace{0.2cm}

		\item[c]
		$\bigcup_{[x] \in A / \sim}[x] = A$

		\vspace{0.2cm}
		\fbox{Proof} ( )
		\vspace{0.2cm}

		$\qed$.

		\vspace{0.5cm}

	\end{enumerate}

%%%%%%%%%%%%%%%%%%%%%%%%%%%%%%%%%%%%%%%%%%%%%%%

\end{document}